\documentclass[12pt,a4paper]{ctexart}
\usepackage{geometry}
\geometry{left=2.5cm,right=2.5cm,top=2.5cm,bottom=2.5cm}
\usepackage{listings}
\usepackage{xcolor}
\usepackage{booktabs}
\usepackage{amsmath}
\usepackage{graphicx}
\usepackage{tcolorbox}
\tcbuselibrary{skins,breakable}
\usepackage[colorlinks,linkcolor=blue!70!black,urlcolor=blue!60!black]{hyperref}

\definecolor{themecolor}{RGB}{30,64,175}
\definecolor{codebg}{RGB}{245,247,250}
\definecolor{codeframe}{RGB}{100,116,139}

\usepackage{titlesec}
\titleformat{\section}{\Large\bfseries\color{themecolor}}{\thesection}{0.8em}{}[\vspace{-0.5em}\rule{\textwidth}{0.4pt}]
\titlespacing*{\section}{0pt}{1.2em}{0.6em}
\titleformat{\subsection}{\large\bfseries\color{themecolor!80}}{\thesubsection}{0.8em}{}
\titlespacing*{\subsection}{0pt}{0.8em}{0.3em}

\lstset{
  basicstyle=\small\ttfamily,
  keywordstyle=\color{blue!80!black}\bfseries,
  commentstyle=\color{green!50!black}\itshape,
  stringstyle=\color{red!70!black},
  numbers=left,numberstyle=\tiny\color{gray},
  backgroundcolor=\color{codebg},
  frame=single,rulecolor=\color{codeframe!40},
  breaklines=true,breakatwhitespace=true,
  language=C++,
  aboveskip=0.8em,belowskip=0.6em
}

\newtcolorbox{webox}{colback=blue!5!white,colframe=themecolor,coltitle=white,fonttitle=\bfseries,title=Web前端展示,boxrule=0.5pt,arc=3pt}

\title{\textbf{编译器设计专题实验报告}\\[0.5em]\Large\color{themecolor}实验六：中间代码生成}
\author{张凌曜 \quad 2234214014 \quad 计算机2306}
\date{}

\begin{document}
\maketitle

\section{实验目的}
写一个中间代码生成器，把源程序翻译成四元式。四元式格式是$(op, arg1, arg2, result)$，表示\texttt{result = arg1 op arg2}，是编译器中间阶段的常用表示。

\section{实验原理}
表达式\texttt{c = a + b * 2}翻译成四元式：

\begin{center}
\begin{tabular}{cccc}
\toprule
op & arg1 & arg2 & result \\
\midrule
* & b & 2 & t1 \\
+ & a & t1 & t2 \\
= & t2 & & c \\
\bottomrule
\end{tabular}
\end{center}

if语句用条件跳转（JZ）实现，while还要在循环体末尾加一个无条件跳转（JMP）回到循环头。

\section{实验过程}

\subsection{表达式翻译}
用递归下降，\texttt{parseExpr}处理加减，\texttt{parseTerm}处理乘除，自底向上翻译。每次遇到二元运算就分配一个临时变量$t_i$，生成一条四元式。乘法在\texttt{parseTerm}里先处理，所以优先级自然是对的。

\subsection{控制流翻译}
if-else需要生成两个标签：L0给else入口，L1给结束位置。条件为假跳L0，if体执行完跳L1，然后L0开始else体，L1结束。

while类似但多了个回跳：L0标记循环头，条件判断后JZ跳出到L1，循环体末尾JMP回L0。

\section{测试结果}

测试程序：

\begin{lstlisting}[language={}]
int main() {
    int a;
    int b;
    int c;
    a = 10;
    b = 20;
    c = a + b * 2;
    if (c > 30) {
        print c;
    } else {
        print a;
    }
    while (a > 0) {
        a = a - 1;
    }
    return 0;
}
\end{lstlisting}

生成的四元式：

\begin{table}[h]
\centering
\begin{tabular}{cllll}
\toprule
\# & op & arg1 & arg2 & result \\
\midrule
000 & FUNC & main & & \\
001--003 & DECL & a,b,c & int & \\
004 & = & 10 & & a \\
005 & = & 20 & & b \\
006 & * & b & 2 & t1 \\
007 & + & a & t1 & t2 \\
008 & = & t2 & & c \\
009 & > & c & 30 & t3 \\
010 & JZ & t3 & & L0 \\
011 & PRINT & c & & \\
012 & JMP & & & L1 \\
013 & LABEL & L0 & & \\
014 & PRINT & a & & \\
015 & LABEL & L1 & & \\
016 & LABEL & L2 & & \\
017 & > & a & 0 & t4 \\
018 & JZ & t4 & & L3 \\
019 & - & a & 1 & t5 \\
020 & = & t5 & & a \\
021 & JMP & & & L2 \\
022 & LABEL & L3 & & \\
023 & RETURN & 0 & & \\
024 & ENDFUNC & main & & \\
\bottomrule
\end{tabular}
\end{table}

看几个关键点：006--008是\texttt{c = a + b * 2}，先算乘法再算加法，优先级没问题。010--015是if-else，JZ跳else，JMP跳过else，标签位置都正确。016--022是while，LABEL标循环头，JZ跳出，JMP回跳。

\begin{center}
\includegraphics[width=\textwidth,height=8cm,keepaspectratio]{image/1.png}
\end{center}

用代码库的\texttt{1.src}也跑了一下，能正确生成四元式。

\begin{center}
\includegraphics[width=\textwidth,height=7cm,keepaspectratio]{image/2.png}
\end{center}

\section{实验总结}
四元式生成的核心就是给每个运算分配临时变量，控制流用标签加跳转处理。这六个实验串起来其实就是一条编译链：DFA定义词法规则 $\to$ 词法分析出Token $\to$ 语法分析建AST $\to$ 语义分析加符号表 $\to$ 中间代码生成四元式，每一步的输出就是下一步的输入。

\begin{webox}
本实验的中间代码生成功能已实现为Web前端页面，与C++版irgen.cpp逻辑一致。

\textbf{主要功能：}源代码输入、四元式中间代码生成（\texttt{(op, arg1, arg2, result)}格式）、符号表展示、Token流展示。支持变量/函数声明、赋值、算术运算（含优先级）、if-else条件跳转、while循环跳转。

\begin{center}
\includegraphics[width=\textwidth,height=8cm,keepaspectratio]{image/web.png}
\end{center}

前端四元式表格按类型着色：LABEL用绿色、JMP/JZ用橙色、FUNC用蓝色，便于直观理解控制流结构。
\end{webox}

\appendix
\section{附录：代码仓库}
本实验的全部源码、报告及Web演示平台已托管至GitHub：

\texttt{https://github.com/yaossss123/compilation-web}

\end{document}
